\documentclass{beamer}
\usepackage[utf8]{inputenc}
\usepackage[T1]{fontenc}
\usepackage{graphicx}
\usepackage{amsmath}
\usepackage{amsfonts}
\usepackage{amssymb}

\usetheme{default}
\usecolortheme{default}

\title{Financial Disclosure in the United States and New Cybersecurity Regulations}
\author{Leonardo Montrucchio}
\date{}

\begin{document}

\begin{frame}
\titlepage
\end{frame}

\begin{frame}{Introduction}
\begin{itemize}
\item In the United States, publicly traded companies are required to periodically submit documents to the Securities and Exchange Commission (SEC) to ensure transparency and investor protection.
\item The main disclosure forms include:
\begin{itemize}
\item Form 10-K (annual report)
\item Form 10-Q (quarterly report)
\item Form 8-K (current report for material events)
\end{itemize}
\end{itemize}
\end{frame}

\begin{frame}{Form 10-K}
\begin{itemize}
\item \textbf{Definition:} Form 10-K is the annual report that publicly traded companies must submit to the SEC, providing a comprehensive overview of their financial performance.
\item \textbf{Content:}
\begin{itemize}
\item Company history and organizational structure.
\item Audited financial statements.
\item Analysis of economic-financial results.
\item Description of business strategy.
\item Discussion of risk factors.
\item Governance information.
\end{itemize}
\item \textbf{Filing deadlines:}
\begin{itemize}
\item Large Accelerated Filers (market capitalization ≥ \$700 million): within 60 days of fiscal year end.
\item Accelerated Filers (market capitalization between \$75 and \$700 million): within 75 days.
\item Non-Accelerated Filers (market capitalization < \$75 million): within 90 days.
\end{itemize}
\end{itemize}
\end{frame}

\begin{frame}{Form 10-K Structure (1/2)}
\textbf{1. Part I (Items 1–4)}
\begin{itemize}
\item Item 1: Business
\item Item 1A: Risk Factors
\item Item 1B: Unresolved Staff Comments
\item Item 1C: Cybersecurity
\item Item 2: Properties
\item Item 3: Legal Proceedings
\item Item 4: Mine Safety Disclosures
\end{itemize}

\textbf{2. Part II (Items 5–9C)}
\begin{itemize}
\item Item 5: Market for Registrant's Common Equity, ...
\item Item 6: Consolidated Financial Data
\item Item 7: Management's Discussion and Analysis (MD\&A)
\item Item 7A: Quantitative and Qualitative Disclosures About Market Risk
\item Item 8: Financial Statements and Supplementary Data
\item Item 9: Changes in and Disagreements With ...
\item Item 9A: Controls and Procedures
\item Item 9B: Other Information
\item Item 9C: Disclosure Regarding Foreign Jurisdictions that Prevent ...
\end{itemize}
\end{frame}

\begin{frame}{Form 10-K Structure (2/2)}
\textbf{3. Part III (Items 10-14)}
\begin{itemize}
\item Item 10: Directors, Executive Officers and Corporate Governance
\item Item 11: Executive Compensation
\item Item 12: Security Ownership of Certain Beneficial Owners ...
\item Item 13: Certain Relationships and Related Transactions, and Director Independence
\item Item 14: Principal Accounting Fees and Services
\end{itemize}

\textbf{4. Part IV (Items 15–16)}
\begin{itemize}
\item Item 15: Exhibits and Financial Statement Schedules
\item Item 16: Form 10–K Summary
\end{itemize}

\vspace{0.5cm}
\footnotesize{Source: Form 10-K General instructions}
\end{frame}

\begin{frame}{Form 10-K Example}
\vspace{2cm}
\begin{center}
\footnotesize{Source: Form 10-K AAPL (2024)}
\end{center}
\end{frame}

\begin{frame}{Form 10-Q}
\begin{itemize}
\item \textbf{Definition:} Form 10-Q is the quarterly report that provides updates on the financial performance of publicly traded companies.
\item \textbf{Content:}
\begin{itemize}
\item Unaudited financial statements.
\item Management analysis of operating results.
\item Updates on risk factors.
\end{itemize}
\item \textbf{Filing deadlines:}
\begin{itemize}
\item Large Accelerated Filers and Accelerated Filers: within 40 days of the end of the fiscal quarter.
\item Non-Accelerated Filers: within 45 days.
\end{itemize}
\end{itemize}
\end{frame}

\begin{frame}{Form 8-K}
\begin{itemize}
\item \textbf{Definition:} Form 8-K is a current report used to communicate material corporate events that occur between periodic filings.
\item \textbf{Examples of events to report:}
\begin{itemize}
\item Changes in corporate leadership.
\item Significant acquisitions or divestitures.
\item Bankruptcies.
\item Material cybersecurity incidents.
\end{itemize}
\item \textbf{Filing timeline:} Companies must file Form 8-K within 4 business days of the event (with possible extension up to 60 days).
\end{itemize}
\end{frame}

\begin{frame}{Form 8-K Structure (1/2)}
\begin{itemize}
\item \textbf{Section 1 – Registrant's Business and Operations}
\begin{itemize}
\item items 1.01 - 1.05 (e.g. Material Definitive Agreements, Bankruptcies, Cybersecurity Incidents)
\end{itemize}
\item \textbf{Section 2 – Financial Information}
\begin{itemize}
\item items 2.01 - 2.06 (e.g. acquisitions and divestitures, operating results, financial obligations)
\end{itemize}
\item \textbf{Section 3 – Securities and Trading Markets}
\begin{itemize}
\item items 3.01 - 3.03 (delisting, sales of unregistered securities, material changes to security holder rights)
\end{itemize}
\item \textbf{Section 4 – Matters Related to Accountants and Financial Statements}
\begin{itemize}
\item items 4.01 - 4.02 (changes in accountant, non-reliance on previously issued financial statements)
\end{itemize}
\end{itemize}
\end{frame}

\begin{frame}{Form 8-K Structure (2/2)}
\begin{itemize}
\item \textbf{Section 5 – Corporate Governance and Management}
\begin{itemize}
\item items 5.01 - 5.08 (e.g. changes in control, resignations or appointments of officers, amendments to articles of incorporation)
\end{itemize}
\item \textbf{Section 6 – Asset-Backed Securities}
\begin{itemize}
\item items 6.01 - 6.06 (e.g. Change of Servicer or Trustee)
\end{itemize}
\item \textbf{Section 7 – Regulation FD Disclosure}
\begin{itemize}
\item item 7.01 (material communications under Reg. FD)
\end{itemize}
\item \textbf{Section 8 – Other Events}
\begin{itemize}
\item item 8.01 (other material events not covered elsewhere)
\end{itemize}
\item \textbf{Section 9 – Financial Statements and Exhibits}
\begin{itemize}
\item item 9.01 (financial statements, pro formas, exhibits and supporting documents)
\end{itemize}
\end{itemize}

\vspace{0.5cm}
\footnotesize{Source: Form 8-K General instructions}
\end{frame}

\begin{frame}{Form 8-K Example}
\vspace{2cm}
\begin{center}
\footnotesize{Source: Form 8-K MSFT (17/01/2024)}
\end{center}
\end{frame}

\begin{frame}{Cybersecurity Disclosure in Form 10-K}
\begin{itemize}
\item \textbf{New SEC Rules (July 2023):} Require publicly traded companies to annually provide information on cybersecurity risk management, strategy and governance in Form 10-K, specifically in Item 1C (Cybersecurity).
\item \textbf{Information to report (Item 106, Regulation S-K):}
\begin{itemize}
\item Risk Management and Strategy: Processes to identify, assess and manage cyber threats, including integration of cybersecurity into business strategy.
\item Cybersecurity Governance: Role of the board of directors and committees in overseeing cyber risks.
\end{itemize}
\item A considerable number of companies also integrate cybersecurity information in Item 1A (Risk Factors), highlighting specific risks associated with cyber threats.
\end{itemize}

\vspace{0.5cm}
\footnotesize{Source: Cybersecurity Disclosure Overview: A Survey of Form 10-K Cybersecurity Disclosures by S\&P 100 Companies (2025)}
\end{frame}

\begin{frame}{S\&P 100 Companies Disclosure Trends}
Item 1C disclosures vary significantly among companies, reflecting each company's specific circumstances. In fact:
\begin{itemize}
\item 66\% of companies divide Item 1C into two main sections: (1) Risk Management and Strategy, and (2) Cybersecurity Governance.
\item 33\% include subsections within one or both main sections.
\item 23\% do not use any specific headings.
\item 11\% adopt headings different from the standard structure, including more than two sections or using alternative titles.
\end{itemize}

\vspace{0.5cm}
\footnotesize{Source: Cybersecurity Disclosure Overview: A Survey of Form 10-K Cybersecurity Disclosures by S\&P 100 Companies (2025)}
\end{frame}

\begin{frame}{Cybersecurity Disclosure in Form 8-K}
\begin{itemize}
\item \textbf{Disclosure obligations:} Since December 18, 2023, companies are required to report "material" cybersecurity incidents in Form 8-K, describing the nature, scope, timing and financial impacts of the event.
\item \textbf{Timeline:} Reporting must occur within 4 business days of determining the materiality of the incident, with possible postponement up to 60 days for national security reasons.
\item \textbf{Materiality impact:} Only 14\% of incidents reported in the first year were considered "material". 28\% of cases reported materiality as "undetermined".
\item \textbf{Alternative items (8.01, 7.01):} When in doubt about materiality, companies have preferred to resort to voluntary disclosures under Item 8.01 or 7.01, in line with SEC advice.
\end{itemize}

\vspace{0.5cm}
\footnotesize{Source: Results from the First Year of Cybersecurity Incident Filings on Form 8-K (2025)}
\end{frame}

\begin{frame}{Reporting Trends and Incident Types}
\begin{itemize}
\item \textbf{Volumes:} In 2024, 55 incidents were reported by 54 companies.
\item \textbf{Timeline:} Average reporting time = 12 days from discovery; 50\% met the 4-day threshold.
\item \textbf{Updates:} About 1/3 of companies filed more than one Form 8-K for the same incident; on average, the first update occurred after 32 days.
\end{itemize}

\vspace{2cm}
\footnotesize{Source: Results from the First Year of Cybersecurity Incident Filings on Form 8-K (2025)}
\end{frame}

\begin{frame}{Our Results}
\begin{itemize}
\item \textbf{Companies analyzed:} 8-K filings from 9,737 US companies published between July 2023 and today were examined.
\item \textbf{Volumes:} 91 Form 8-K filings were identified from 66 different companies, relating to 69 cyber attacks.
\item \textbf{Timeline:} The average reporting time found is approximately 12 days. In 56\% of cases, the 4-day threshold was met.
\item \textbf{Updates:} 18 companies filed more than one Form 8-K, with an average update time of 36 days. 2 companies first declared via item 8.01 and then 1.05, formalizing the materiality of the event.
\item \textbf{Other info:} 1 company received more than one attack.
\end{itemize}
\end{frame}

\begin{frame}{Bibliography}
\footnotesize{
SEC (2024), Form 10-K General Instructions. U.S. Securities and Exchange Commission.

SEC (2024), Form 8-K General Instructions. U.S. Securities and Exchange Commission.

Apple Inc. (2024), Form 10-K AAPL (2024).

Microsoft Corporation (2024), Form 8-K MSFT (17/01/2024).

Cybersecurity Disclosure Overview (2025), A Survey of Form 10-K Cybersecurity Disclosures by S\&P 100 Companies. Harvard Law School Forum on Corporate Governance.

Kayne McGladrey (2025), Results from the First Year of Cybersecurity Incident Filings on Form 8-K.
}
\end{frame}

\end{document}